\section{\texorpdfstring{\texttt{structure} versus \texttt{class}: why this book delays type classes}{structure versus class: why this book delays type classes}}\label{sec:05-rigor-check:01-structure-vs-class:1}

Every algebraic definition in this book --- \texttt{Group}, \texttt{Ring}, \texttt{Module} --- is written as a plain \texttt{structure}. Mathlib, the real Lean library, defines all of these using \texttt{class} instead:

\begin{lstlisting}
-- this book's style
structure Group (G : Type) where
  op : G → G → G
  id : G
  inv : G → G
  assoc : ∀ a b c : G, op (op a b) c = op a (op b c)
  id_left : ∀ a : G, op id a = a
  -- ...

-- Mathlib's style (schematically — not the literal Mathlib source)
class Group (G : Type) where
  op : G → G → G
  id : G
  inv : G → G
  assoc : ∀ a b c : G, op (op a b) c = op a (op b c)
  id_left : ∀ a : G, op id a = a
  -- ...
\end{lstlisting}

Syntactically these look almost identical. \texttt{class} is very nearly a shorthand for \texttt{structure} plus one extra mechanism. Understanding exactly what that mechanism is, and why this book avoids it until the intuition \texttt{class} would otherwise hide has been built, is the point of this section.

\subsection{\texorpdfstring{What \texttt{class} actually adds: instance search}{What class actually adds: instance search}}\label{sec:05-rigor-check:01-structure-vs-class:2}

A \texttt{class} is a \texttt{structure} whose \emph{inhabitants} Lean is willing to search for automatically, through a process called \textbf{typeclass resolution} (or ``instance search''). When you write \texttt{class Group (G : Type)}, and then register a particular group structure with the keyword \texttt{instance} instead of \texttt{def}:

\begin{lstlisting}
instance : Group Int where
  op := fun a b => a + b
  id := 0
  inv := fun a => -a
  -- ...
\end{lstlisting}

Lean records \texttt{Int}'s group structure in a global table. From then on, any function that takes an \emph{implicit} \texttt{[Group G]} argument (square brackets, not curly braces --- this is a third kind of argument, an \textbf{instance argument}) has it filled in automatically whenever \texttt{G} is unified with \texttt{Int}, with no explicit term needed at the call site:

\begin{lstlisting}
def opTwice [Group G] (x : G) : G :=
  Group.op x x   -- the `Group G` instance is found automatically

#eval opTwice (3 : Int)   -- Lean infers the `Group Int` instance silently
\end{lstlisting}

Compare this to the style used throughout this book, where every theorem about groups takes an \emph{explicit} term \texttt{Grp : Group G} as a genuine function argument:

\begin{lstlisting}
def opTwice (Grp : Group G) (x : G) : G :=
  Grp.op x x

#eval opTwice intGroup 3   -- the specific Group term must be named
\end{lstlisting}

Mathematically, the difference is this: \texttt{class} + \texttt{instance} implements the same convention as ordinary mathematical prose, in which one writes ``let \(G\) be a group'' once and then simply writes \(a \cdot b\), \(e\), \(a^{-1}\) thereafter. The \emph{specific} group structure is left implicit, tracked by context, and never re-named at each use. Lean's typeclass mechanism automates exactly that convention. A \texttt{structure}-only approach (this book's style) requires the structure to always be carried around explicitly, so the group operation is never left unstated.

\subsection{\texorpdfstring{Why this book deliberately avoids \texttt{class}, for now}{Why this book deliberately avoids class, for now}}\label{sec:05-rigor-check:01-structure-vs-class:3}

\texttt{class}'s automation is exactly what would make it harder to see, the first time through, what is actually happening. Every \texttt{Grp.op}, every \texttt{Grp.assoc} in this book's proofs is a visible, traceable term. One can always ask ``where did this fact come from?'' and point at the specific argument it is a field of. With \texttt{class}, that same information is present but resolved silently by the elaborator --- convenient once trusted, but difficult to debug the first time instance resolution picks the ``wrong'' (or an unexpected) instance, or fails to find one at all. Reading the \emph{explicit} version first, this book's approach, builds the mental model that makes debugging \emph{silent} instance-resolution failures manageable later, since the data being threaded through remains known even once no longer visible in the source text.

\subsection{The bridge to Mathlib}\label{sec:05-rigor-check:01-structure-vs-class:4}

Nothing about the \emph{mathematical content} changes between the two styles. \texttt{Group}'s axioms are the axioms, in either encoding. The type-class version also usually bundles in the group's \texttt{*}, \texttt{⁻¹}, \texttt{1} as genuine \emph{notation} (through further, separate mechanisms --- \texttt{Mul}, \texttt{Inv}, \texttt{One} type classes that \texttt{Group} extends), so that Mathlib code reads as \texttt{a * b} instead of \texttt{Grp.op a b}. Chapter 13's suggested projects include redoing this book's \texttt{Group}/\texttt{Ring} as type classes once the reader is comfortable doing so. At that point, everything in this section is the vocabulary needed.

\subsection{\texorpdfstring{A note on \texttt{.toGroup}, \texttt{.toPoint}, and coercion}{A note on .toGroup, .toPoint, and coercion}}\label{sec:05-rigor-check:01-structure-vs-class:5}

Whenever a \texttt{structure} extends another (Chapter 2's \texttt{Point3D extends Point}, this book's \texttt{CommGroup extends Group}), Lean generates a field named \texttt{.toX} that projects back to the parent structure. \texttt{cg.toGroup} recovers the underlying \texttt{Group} from a \texttt{CommGroup cg}. Lean also permits \emph{dropping} the \texttt{.toGroup} and writing \texttt{cg.op} directly, silently inserting the projection. This automatic insertion of a ``translate from one type to another'' function is an instance of \textbf{coercion}, the general mechanism by which Lean lets a term of one type stand in for a term of a related type without the conversion being written explicitly. It is the same general idea, though a much simpler case, as coercing \texttt{Nat} values into \texttt{Int}. The literal projection differs, but ``insert a conversion function automatically so notation reads naturally'' is the shared mechanism.

\begin{quote}
Read more: TPiL's chapters on ``Structures and Records'' and ``Type Classes'' cover both topics in this section directly, including coercions, in more depth and with additional worked examples than this book provides.
\end{quote}
